%% ma: L10-B13-0067
%% lop: 10
%% chuong: 5
%% bai: 13
%% dang: 10.13.13
%% loai: TL
%% muc: TH
%% dap_an: "-"
%% nguon: "chuyen-de-cac-so-dac-trung-cua-mau-so-lieu-khong-ghep-nhom-toan-10-chuong-trinh-moi.pdf (GV. Phan Nhật Linh) (P03), Bài 02 Dạng 2 Bài tập tự luận 4"
%% kien_thuc: "Lớp 10 Bài 13 (số trung bình, trung vị, tứ phân vị, mốt)"
%% trang_thai: chinh-thuc
%% ghi_chu: "Bài tập tự luận 8 ý a), b) trùng hoàn toàn bài này, đã bỏ hai ý đó. Tự duyệt theo yêu cầu giáo viên 22/10/2026."
\begin{ex}
Tìm số trung bình, các tứ phân vị và mốt của các mẫu số liệu sau:

a) Số điểm mà năm vận động viên bóng rổ ghi được trong một trận đấu: $9$; $8$; $15$; $8$; $20$.

b) Giá của một số loại giày (đơn vị: nghìn đồng): $350$; $300$; $650$; $300$; $450$; $500$; $300$; $250$.
\loigiai{a) Sắp xếp: $8$; $8$; $9$; $15$; $20$. $\overline x=\dfrac{60}5=12$; $Q_2=9$; $Q_1=\dfrac{8+8}2=8$; $Q_3=\dfrac{15+20}2=17{,}5$; mốt $8$.

b) Sắp xếp: $250$; $300$; $300$; $300$; $350$; $450$; $500$; $650$. $\overline x=\dfrac{3100}8=387{,}5$; $Q_2=\dfrac{300+350}2=325$; $Q_1=\dfrac{300+300}2=300$; $Q_3=\dfrac{450+500}2=475$; mốt $300$.}
\end{ex}
